\subsection{Perkalian Skalar pada Matriks}

Ingat bahwa kita menggunakan istilah {\em skalar \em} untuk menyebut bilangan. Oleh karena itu, {\em perkalian skalar pada matriks \em} adalah perkalian suatu matriks dengan sebuah bilangan.  
Untuk mengilustrasikan konsep ini, perhatikan contoh berikut, di mana sebuah
matriks dikalikan dengan skalar $3$.
\begin{equation*}
3\leftB
\begin{array}{rrrr}
1 & 2 & 3 & 4 \\
5 & 2 & 8 & 7 \\
6 & -9 & 1 & 2
\end{array}
\rightB = \leftB
\begin{array}{rrrr}
3 & 6 & 9 & 12 \\
15 & 6 & 24 & 21 \\
18 & -27 & 3 & 6
\end{array}
\rightB 
\end{equation*}

Matriks baru diperoleh dengan mengalikan setiap entri matriks semula
dengan skalar yang diberikan. 

Definisi formal perkalian skalar adalah sebagai berikut.

\begin{definition}{Perkalian Skalar pada Matriks}\label{def:scalarmultofmatrices}
Jika $A=\leftB a_{ij}\rightB $ dan $k$ adalah sebuah skalar,
maka $kA=\leftB ka_{ij}\rightB .$ \index{matriks!perkalian skalar}
\end{definition}

Perhatikan contoh berikut.

\begin{example}{Pengaruh Perkalian dengan Skalar}\label{exa:effectofscalarmult}
Tentukan hasil perkalian matriks $A$ berikut dengan $7$.
\begin{equation*}
A=\leftB
\begin{array}{rr}
2 & 0 \\
1 & -4
\end{array}
\rightB
\end{equation*}
\end{example}

\begin{solution}
Berdasarkan Definisi \ref{def:scalarmultofmatrices}, kita mengalikan setiap elemen $A$ dengan $7$.
Oleh karena itu,
\begin{equation*}
7A = 
7\leftB
\begin{array}{rr}
2 & 0 \\
1 & -4
\end{array}
\rightB =
\leftB
\begin{array}{rr}
7(2) & 7(0) \\
7(1) & 7(-4)
\end{array}
\rightB =
\leftB
\begin{array}{rr}
14 & 0 \\
7 & -28
\end{array}
\rightB
\end{equation*}
\end{solution}

Seperti halnya penjumlahan matriks, ada beberapa sifat perkalian skalar yang berlaku.

\begin{proposition}{Sifat-Sifat Perkalian Skalar}\label{prop:propertiesscalarmult}
Misalkan $A, B$ adalah matriks, dan $k, p$ adalah skalar. Maka sifat-sifat berikut \index{matriks!sifat perkalian skalar} berlaku.
\begin{itemize}
\item Hukum Distributif terhadap Penjumlahan Matriks
\begin{equation*}
k \left( A+B\right) =k A+ kB  
\end{equation*}

\item Hukum Distributif terhadap Penjumlahan Skalar
\begin{equation*}
\left( k +p \right) A= k A+p A
\end{equation*}

\item Hukum Asosiatif untuk Perkalian Skalar
\begin{equation*}
k \left( p A\right) = \left( k p \right) A 
\end{equation*}

\item Aturan Perkalian dengan $1$
\begin{equation*}
1A=A  
\end{equation*}
\end{itemize}

\end{proposition}

Bukti proposisi ini serupa dengan bukti Proposisi \ref{prop:propertiesofaddition} dan diserahkan sebagai latihan kepada pembaca.
